% 2013 四川卷文科第 7 题：四个频率分布直方图选项 (A)—(D)。
% 组距 5 时 频率/组距 = 频数/100：(A) 频数 1,1,4,2,4,3,3,2（正确），
% (B) 1,2,3,3,5,2,3,1；(C)(D) 误用组距 10。
\newcommand{\freqpanel}[5]{%
  \begin{scope}[shift={#1}, x=0.142cm, y=58cm]
    \draw[->] (0,0) -- (44,0) node[above left, inner sep=3pt] {\small 人数};
    \draw[->] (0,0) -- (0,{#2*1.36});
    \node[above, inner sep=1pt] at (0,{#2*1.42})
      {$\dfrac{\text{频率}}{\text{组距}}$};
    % 纵轴刻度
    \foreach \v in {#4} {
      \draw (-0.9,\v) -- (0,\v);
      \node[left, inner sep=1.5pt] at (-1.0,\v) {\v};
    }
    % 横轴刻度
    \foreach \t in {#3} {
      \draw (\t,-0.0016) -- (\t,0.0016);
      \node[below, inner sep=1.5pt] at (\t,-0.0019) {\t};
    }
    % 柱形
    \foreach \h [count=\i from 1] in {#5} {
      \draw ({(\i-1)*\barw},0) rectangle ({\i*\barw},\h);
    }
  \end{scope}%
}
\begin{tikzpicture}[font=\scriptsize, line width=0.45pt, >=Stealth]
  \def\barw{5}
  \freqpanel{(0,0)}{0.04}{0,5,...,40}{0.01,0.02,0.03,0.04}
    {0.01,0.01,0.04,0.02,0.04,0.03,0.03,0.02}
  \node at (-0.55,-1.15) {(A)};
  \freqpanel{(7.4,0)}{0.05}{0,5,...,40}{0.01,0.02,0.03,0.04,0.05}
    {0.01,0.02,0.03,0.03,0.05,0.02,0.03,0.01}
  \node at (6.85,-1.15) {(B)};
  \def\barw{10}
  \freqpanel{(0,-5.6)}{0.04}{0,10,20,30,40}{0.01,0.02,0.03,0.04}
    {0.01,0.03,0.035,0.025}
  \node at (-0.55,-6.75) {(C)};
  \freqpanel{(7.4,-5.6)}{0.04}{0,10,20,30,40}{0.01,0.02,0.03,0.04}
    {0.015,0.03,0.035,0.02}
  \node at (6.85,-6.75) {(D)};
\end{tikzpicture}
